\begin{exercise}[RG-improved fermion branch point]{I-CH09-016}
Let $\Delta=p^2+m^2$ denote the Euclidean distance from the fermion mass shell.
Near $\Delta=0$, write the renormalized propagator as
\[
  S_R(p)=\frac{\slashed p+\ii m}{\Delta}\,F(\Delta/\mu^2,\alpha,\kappa).
\]
Use the connected two-point RG equation
\[
  \mathcal D S_R=-2\gamma_F S_R,
  \qquad
  \gamma_F^{(1)}=-\frac{\kappa\alpha}{4\pi},
  \qquad
  \beta(\alpha)=\frac{2\alpha^2}{3\pi}+O(\alpha^3),
\]
with Banks's convention for $\gamma_F$.  Solve the characteristic equation,
first keeping $\alpha$ and $\kappa$ fixed at leading-log order and then writing
the result with their running values.  Show that the pole becomes a
gauge-dependent noninteger power of $\Delta$ and therefore a branch point.
Compute the discontinuity phase after continuation across the physical cut.
Check that the leading branch exponent vanishes in Landau gauge, and explain
why this singularity cannot by itself define a gauge-invariant charged-particle
$S$-matrix state.
\end{exercise}
